% Part: lambda-calculus
% Chapter: introduction
% Section: syntax

\documentclass[../../../include/open-logic-section]{subfiles}

\begin{document}

\olfileid[id]{lam}{int}{syn}
\olsection{Sintaksis Kalkulus Lambda}

Kita mulai dengan suatu barisan variabel $x$, $y$, $z$,~\dots dan beberapa
simbol konstanta $a$, $b$, $c$,~\dots. Himpunan suku didefinisikan secara
induktif sebagai berikut:
\begin{enumerate}
\item Setiap variabel adalah suku.
\item Setiap konstanta adalah suku.
\item Jika $M$ dan $N$ adalah suku, maka $(MN)$ juga suku.
\item Jika $M$ adalah suku dan $x$ adalah variabel, maka
  $(\lambd[x][M])$ adalah suku.
\end{enumerate}
Suku berbentuk $(MN)$ disebut \emph{aplikasi}, sedangkan suku berbentuk
$(\lambd[x][M])$ disebut \emph{abstraksi}.

Sistem tanpa konstanta sama sekali disebut kalkulus lambda \emph{murni}. Kita
terutama akan bekerja dalam kalkulus-$\lambd$ murni, sehingga semua huruf kecil
menyatakan variabel. Kita menggunakan huruf besar ($M$, $N$, dan seterusnya)
untuk menyatakan suku-suku kalkulus-$\lambd$.

Kita akan mengikuti beberapa konvensi notasi:

\begin{conv}
\begin{enumerate}
\item Apabila tanda kurung dihilangkan, aplikasi dikelompokkan dari kiri ke
  kanan. Sebagai contoh, jika $M$, $N$, $P$, dan $Q$ adalah suku, maka $MNPQ$
  menyingkat $(((MN)P)Q)$.
\item Sekali lagi, apabila tanda kurung dihilangkan, abstraksi lambda diberi
  cakupan seluas mungkin. Sebagai contoh, $\lambd[x][MNP]$ dibaca sebagai
  $(\lambd[x][((MN)P)])$.
\item Satu lambda dapat digunakan untuk mengabstraksikan beberapa variabel.
  Sebagai contoh, $\lambd[xyz][M]$ menyingkat
  $\lambd[x][\lambd[y][\lambd[z][M]]]$.
\end{enumerate}
\end{conv}

Sebagai contoh,
\[
\lambd[xy][xxyx \lambd[z][xz]]
\]
menyingkat
\[
\lambd[x][\lambd[y][((((xx)y)x)(\lambd[z][(xz)]))]].
\]
Konvensi-konvensi ini perlu dihafalkan. Pada awalnya konvensi tersebut akan
membingungkan, tetapi Anda akan terbiasa; setelah beberapa waktu, konvensi itu
akan jauh kurang membingungkan daripada harus menghadapi belantara tanda
kurung.

Dua suku yang hanya berbeda dalam nama variabel terikat disebut
$\alpha$-ekuivalen; misalnya, $\lambd[x][x]$ dan $\lambd[y][y]$. Akan lebih
praktis jika keduanya dipandang sebagai suku yang ``sama''; dengan kata lain,
ketika kita mengatakan bahwa $M$ dan $N$ sama, yang dimaksud juga mencakup
``hingga penggantian nama variabel terikat.'' Variabel yang berada dalam
cakupan suatu $\lambd$ disebut ``terikat'', sedangkan variabel lainnya disebut
``bebas.'' Tidak ada variabel bebas dalam contoh sebelumnya; tetapi dalam
\[
(\lambd[z][yz])x
\]
$y$ dan $x$ bebas, sedangkan $z$ terikat.

\end{document}
